\subsection{API klientów}

Rejestr klientów przechowuje dane kupujących, nie konta z hasłami.
Ścieżka bazowa to \texttt{/clients}. Kontroler zwraca bezpośrednio
\hyperref[schema:client]{obiekty Client}, analogicznie do katalogu ofert.

\subsubsection{Endpointy API}

\begin{table}[H]
  \centering
  \renewcommand{\arraystretch}{1.3}
  \begin{tabularx}{\textwidth}{|l|l|X|}
    \hline
    \textbf{Endpoint} & \textbf{Metoda} & \textbf{Opis} \\
    \hline
    \texttt{/clients} & \texttt{POST} & Tworzy klienta na podstawie pól \texttt{name} i \texttt{email}. UUID nadaje serwer. \\
    \hline
    \texttt{/clients} & \texttt{GET} & Zwraca listę klientów. \\
    \hline
    \texttt{/clients/\{id\}} & \texttt{GET} & Zwraca klienta o podanym UUID. \\
    \hline
    \texttt{/clients/\{id\}} & \texttt{PUT} & Zastępuje oba pola: \texttt{name} i \texttt{email}. UUID nie zmienia się. \\
    \hline
    \texttt{/clients/\{id\}} & \texttt{DELETE} & Usuwa klienta bez rezerwacji i zwraca usunięty obiekt. \\
    \hline
  \end{tabularx}
  \caption{Endpointy API klientów (\texttt{/clients})}
  \label{tab:client-endpoints}
\end{table}

\begin{figure}[H]
\begin{verbatim}
{
  "name": "Jan Kowalski",
  "email": "jan@example.org"
}
\end{verbatim}
\caption{Żądanie utworzenia lub aktualizacji klienta}
\label{schema:client-request}
\end{figure}

Po usunięciu otaczających białych znaków \texttt{name} musi mieć od 1 do 255 znaków,
a \texttt{email} nie więcej niż 254 znaki i poprawną podstawową składnię adresu.
Nie jest sprawdzana dostarczalność poczty. E-mail nie jest unikalnym kluczem.
Rezerwacja przechowuje UUID klienta, nie kopię jego danych kontaktowych.

\subsubsection{Kody odpowiedzi API klientów}

\begin{table}[H]
  \centering
  \renewcommand{\arraystretch}{1.3}
  \begin{tabularx}{\textwidth}{|l|l|X|}
    \hline
    \textbf{Endpoint} & \textbf{Kod odpowiedzi} & \textbf{Opis} \\
    \hline
    Wszystkie & 200 OK & Operacja zakończona poprawnie; odpowiedzią jest klient lub lista klientów. \\
    \hline
    \texttt{POST}, \texttt{PUT} & 400 Bad Request & Brak wymaganych pól lub niepoprawne dane. \\
    \hline
    Operacje na \texttt{/\{id\}} & 404 Not Found & Klient nie istnieje, także dla niepoprawnego identyfikatora w ścieżce. \\
    \hline
    \texttt{DELETE /clients/\{id\}} & 409 Conflict & Klient ma rezerwacje i nie może zostać usunięty. \\
    \hline
  \end{tabularx}
  \caption{Kody odpowiedzi endpointów klientów}
  \label{tab:client-response-codes}
\end{table}

\subsection{API rezerwacji}

Ścieżka bazowa to \texttt{/reservations}. Kontroler zwraca bezpośrednio
\hyperref[schema:reservation]{obiekty Reservation}, a obsługa wyjątków jest
zdefiniowana w kontrolerze, tak jak w katalogu ofert.

\subsubsection{Endpointy API}

\begin{table}[H]
  \centering
  \renewcommand{\arraystretch}{1.3}
  \begin{tabularx}{\textwidth}{|l|l|X|}
    \hline
    \textbf{Endpoint} & \textbf{Metoda} & \textbf{Opis} \\
    \hline
    \texttt{/reservations} & \texttt{POST} & Tworzy koszyk na podstawie \texttt{clientId} oraz niepustej listy \texttt{items}. \\
    \hline
    \texttt{/reservations/\{id\}} & \texttt{GET} & Zwraca rezerwację o podanym UUID. \\
    \hline
    \texttt{/reservations} & \texttt{GET} & Zwraca rezerwacje klienta wskazanego wymaganym parametrem \texttt{clientId}. \\
    \hline
  \end{tabularx}
  \caption{Endpointy API rezerwacji (\texttt{/reservations})}
  \label{tab:reservation-endpoints}
\end{table}

\begin{figure}[H]
\begin{verbatim}
{
  "clientId": "11111111-1111-4111-8111-111111111111",
  "items": [
    {
      "offerId": "22222222-2222-4222-8222-222222222222",
      "quantity": 2
    }
  ]
}
\end{verbatim}
\caption{Żądanie utworzenia rezerwacji}
\label{schema:create-reservation-request}
\end{figure}

Podane UUID-y muszą wskazywać istniejącego klienta i oferty. Serwer nadaje UUID
rezerwacji i pozycji, znaczniki czasu oraz pobiera zrzuty ofert z katalogu.
Przekazane dodatkowo ceny lub zrzuty nie nadpisują wartości katalogowych.
Statusy i dane uczestników nie są przyjmowane jako pola modelu ani zapisywane.

Pole \texttt{quantity} jest dodatnią liczbą całkowitą JSON z zakresu typu Java
\texttt{int}. Wartości ułamkowe i tekstowe są odrzucane, nie konwertowane.
Suma ilości dla powtarzającej się oferty również musi mieścić się w tym zakresie.
Powtarzające się oferty zachowują osobne pozycje koszyka, ale mają wspólny zrzut
i jedną zagregowaną operację odjęcia zapasu.

\subsubsection{Kody odpowiedzi API rezerwacji}

\begin{table}[H]
  \centering
  \renewcommand{\arraystretch}{1.3}
  \begin{tabularx}{\textwidth}{|l|l|X|}
    \hline
    \textbf{Endpoint} & \textbf{Kod odpowiedzi} & \textbf{Opis} \\
    \hline
    Wszystkie & 200 OK & Utworzona lub odczytana rezerwacja, albo lista rezerwacji klienta. \\
    \hline
    \texttt{POST /reservations} & 400 Bad Request & Brak pozycji, błędny UUID, niepoprawna ilość lub przekroczenie zakresu sumy ilości. \\
    \hline
    \texttt{GET /reservations} & 400 Bad Request & Brak parametru \texttt{clientId}. \\
    \hline
    Wszystkie & 404 Not Found & Wskazany klient, oferta lub rezerwacja nie istnieje. \\
    \hline
    \texttt{POST /reservations} & 409 Conflict & Niedostateczny zapas lub konflikt wersji oferty. \\
    \hline
  \end{tabularx}
  \caption{Kody odpowiedzi endpointów rezerwacji}
  \label{tab:reservation-response-codes}
\end{table}

\subsection{Walidacja, spójność i zapis koszyka}

\begin{itemize}
  \item \textbf{Granica modułu ofert}: \texttt{ReservationService} korzysta z
    \texttt{OfferChangeService}, nie z repozytorium ofert.
    \texttt{OfferChangeService.reserveQuantity} sprawdza w kodzie Java dodatniość
    ilości oraz dostępny zapas. Przy braku zapasu zgłasza konflikt przed zapisem.
    Nowa dostępność jest obliczana w usłudze i zapisywana istniejącą metodą
    aktualizacji z kontrolą wersji.
  \item \textbf{Współbieżność}: aktualizacja porównuje wersję oferty i zwiększa ją.
    Warunek wersji zapobiega nadpisaniu nowszej zmiany; nie zastępuje walidacji
    ilości w usłudze. Istniejące ograniczenie schematu dla nieujemnej dostępności
    pozostaje dodatkową ochroną, nie mechanizmem walidacji żądań.
  \item \textbf{Transakcja}: utworzenie rezerwacji, odjęcie wszystkich ilości i zapis
    JSON pozycji są jedną transakcją. Dowolny błąd cofa cały koszyk i zmiany zapasu.
    Pierwsza operacja, \texttt{INSERT} rezerwacji dla istniejącego klienta,
    uzyskuje blokadę zapisu SQLite przed odczytem ofert. Pusty koszyk roboczy
    nigdy nie jest zatwierdzany.
  \item \textbf{Historia}: \texttt{offer\_items\_data} zawiera pełne pozycje ze
    zrzutami ofert. Późniejsze zmiany lub usunięcie oferty nie zmieniają rezerwacji.
    \texttt{offerSnapshot.quantity} to zapas przed zakupem, a ilość pozycji
    oznacza liczbę rezerwowanych sztuk. Ceny zachowują tekstową reprezentację
    \texttt{Money} z dwoma miejscami po przecinku.
  \item \textbf{Czas}: w bazie znaczniki czasu są liczbami milisekund od epoki Unix,
    w API wartościami ISO-8601 w UTC. Przy utworzeniu wszystkie znaczniki są równe.
    Wersja rezerwacji wynosi 1 i nie jest eksponowana przez API.
  \item \textbf{Repozytoria}: szablony zapytań SQL są stałymi w klasach repozytoriów,
    analogicznie do istniejącego repozytorium ofert.
\end{itemize}

\subsection{Konfiguracja i ograniczenia}

\begin{itemize}
  \item \textbf{Klucze obce}: \texttt{connection-init-sql} w konfiguracji Hikari
    wykonuje \texttt{PRAGMA foreign\_keys = ON} na każdym nowym połączeniu.
    Zapewnia to egzekwowanie relacji, w szczególności rezerwacja--klient.
    Kolejność usuwania tabel w skrypcie inicjalizacyjnym uwzględnia te zależności.
  \item \textbf{Oczekiwanie na blokadę}: właściwość sterownika SQLite
    \texttt{busy\_timeout=10000} pozwala czekać do 10 sekund na zwolnienie blokady.
    Nie jest to ponawianie żądania HTTP ani zwiększenie liczby równoległych
    zapisów SQLite.
  \item \textbf{Porządek konfiguracji}: usunięto duplikaty ustawień źródła danych
    i inicjalizacji SQL spod sekcji \texttt{server}, która nie konfiguruje tych
    mechanizmów Spring Boot. Właściwe ustawienia pozostają pod \texttt{spring}.
    Port 8081, plik \texttt{offers.db} i tryb inicjalizacji nie zostały zmienione.
  \item \textbf{Inicjalizacja}: domyślne \texttt{spring.sql.init.mode=always}
    odtwarza schemat i usuwa wszystkie dane przy starcie, także klientów
    i rezerwacje. Po pierwszej inicjalizacji można zachować bazę przez ustawienie
    \texttt{SPRING\_SQL\_INIT\_MODE=never}. Migracje schematu pozostają poza zakresem.
  \item \textbf{Zakres}: statusy, dane i reguły uczestników, płatności,
    anulowanie i zwalnianie zapasu, wygasanie, konkretne miejsca oraz oferty
    zależne nie są zaimplementowane. Rezerwacje zajmują zapas bezterminowo,
    ponowiony \texttt{POST} tworzy kolejny koszyk, a listy nie mają paginacji.
  \item \textbf{Dostęp}: projekt nie ma uwierzytelniania ani autoryzacji.
    Endpointów udostępniających dane klientów nie należy wystawiać publicznie.
\end{itemize}
